% Add "answers" option to print the answers
\documentclass[addpoints]{exam}

\usepackage[utf8]{inputenc}
\usepackage{calc}
\usepackage[dvipsnames]{xcolor}
\usepackage{ifthen}
\usepackage{tcolorbox}

% \makebox[\textwidth]{Name:\enspace\hrulefill}
%
% \vspace{5mm}
%\makebox[\textwidth]{Instructor’s name:\enspace\hrulefill}

\definecolor{MyLightGray}{gray}{0.60}

\newcommand{\truefalse}
{		\begin{checkboxes}
			 \choice True
			 \choice False. Explanation:
		\end{checkboxes}
		\vspace{10pt}
}


% Check if 'answers' is an option
\newcommand{\ifanswers}[2]{%
  \ifthenelse{\boolean{exam@answers}}{#1}{#2}%
}



\newcommand{\answerspace}[1]
{
	\ifprintanswers
	\else
		\setlength{\mylength}{10pt}%
		% \vspace{#1\mylength}
		% Make the grey box
		\par\nobreak\noindent
		\makebox[\linewidth][l]{%
			% set inner padding to zero, so only the frame shows
			{\setlength{\fboxsep}{0pt}%
			% draw a color frame, white (or transparent) inside
			\fcolorbox{MyLightGray}{white}{%
			% the inner box has zero width content but fixed height
			\parbox[t][#1\mylength][c]{%
				% subtract twice the frame‐thickness from the full width
				\dimexpr\linewidth-2\fboxrule\relax
			}{\strut}%
			}%
			}%
		}%
		\par
		\vspace{3\mylength}

  	\fi
}
\newlength{\mylength}

\newcommand{\code}[1]{{\texttt{#1}}}


%\usepackage{listings}
\usepackage{color}
\usepackage{lastpage}
\usepackage{verbatim}
\usepackage{fancyvrb}
\usepackage{graphicx}
\usepackage{float}
\usepackage{minted}
\usepackage{multicol}

\newcommand{\includecode}[2][c]{\lstinputlisting[caption=, escapechar=, language=#1]{#2}}
\RecustomVerbatimCommand{\VerbatimInput}{VerbatimInput}{}

\newcommand{\Code}[1]{\mintinline{python}{#1}}

\renewcommand{\solutiontitle}{\noindent\textbf{Solution sketch: }}

\definecolor{SolutionColor}{rgb}{0.8,0.9,1}

\cfoot{Page \thepage\ of \pageref{LastPage}}

\rhead{Initials: \rule{20mm}{0.4pt}}


%%%%%%%%%%%%%%%%%%%%%%%%%%%%%%%%%%%%%%%%
\begin{document}
\fvset{numbers=left,numbersep=3pt}
\thispagestyle{empty}

\begin{center}
	{\LARGE CS 408 Fall 2026 -- Mutation-Based Fuzzing Quiz \\}
\fbox{\fbox{\parbox{5.5in}{\centering
	\begin{itemize}
        \item For this quiz, think of a solution to each question on
          your own, then discuss with your friends and neighbors, and
          then write down the answers you collaboratively came up
          with.

        \item Full points will be rewarded for \textbf{any} (reasonable)
          attempt.
        \end{itemize}
	}}}
\end{center}

\vspace{5mm}
\makebox[\textwidth]{\textbf{Name (please print)}: \rule{100mm}{0.8pt} }

\vspace{10mm}
\makebox[\textwidth]{\textbf{PUID}: \rule{100mm}{0.8pt} \hspace{35mm}}

\begin{questions}

\question A mutation-based fuzzer is testing a web browser. Its seed is
  a single valid URL:
  \begin{center}
    \Code{http://cs.purdue.edu/index.html}
  \end{center}
  Recall the three mutation operators: \textbf{delete} a character at a
  random position, \textbf{insert} a random printable character at a
  random position, and \textbf{flip} one bit of the character at a
  random position.

\begin{parts}

\part What is the name of one classmate you discussed your answers
  with?
  \answerspace{5} %
  \begin{solution}
    No wrong answers here!
  \end{solution}

\part Give one single mutation that leaves the URL \textbf{still
  valid}, and one that \textbf{breaks it immediately}. Be specific about
  the position and the operator. %
  \answerspace{24} %
  \begin{solution} %
    Many correct answers. For example:

    \textbf{Still valid:} delete the \Code{'c'} of \Code{cs},
    giving \Code{http://s.purdue.edu/index.html}. Or flip bit 5 of the
    leading \Code{'h'}, giving
    \Code{Http://cs.purdue.edu/index.html} --- schemes are
    case-insensitive.

    \textbf{Breaks it:} delete the \Code{':'}, giving
    \Code{http//cs.purdue.edu/index.html}, which has no scheme
    separator. Or delete any character of \Code{http}.

    For the record: of the 29 possible single deletions, 24 leave a
    valid URL, and of the 217 possible single bit flips, 171 do. Most
    single mutations \emph{do} preserve validity --- which is the whole
    bet behind mutating a seed rather than generating from scratch.
  \end{solution}

\part Now suppose the fuzzer applies \textbf{ten} mutations in a row
  rather than one. Would you expect the result to still be a valid URL?
  What does your answer say about how a fuzzer should spend its budget? %
  \answerspace{24} %
  \begin{solution} %
    Much less likely. Each mutation is roughly an 80\% chance of
    staying valid, so ten in a row is around $0.8^{10} \approx 11\%$
    --- and in practice worse, since damage accumulates in the parts
    of the URL that matter.

    So a fuzzer should not simply mutate a seed many times. It should
    mutate \emph{once}, keep the result if it is worth keeping, and
    mutate \emph{that}. Which raises the question of part (c).
  \end{solution}

% \newpage

% \part Of the mutations that keep the URL valid, most are not
%   interesting: they explore the same code the seed already did. What
%   could the fuzzer \textbf{measure} to tell an interesting mutant from
%   a boring one? How would it use that measurement? %
%   \answerspace{20} %
%   \begin{solution} %
%     Measure \textbf{coverage}: run the program on the mutant and record
%     which lines or branches it executed. If the mutant reaches
%     something none of the kept inputs reached, it is worth keeping;
%     otherwise throw it away.

%     This is the same measurement we used in week 3 to \emph{grade} a
%     test suite. Here it is used to \emph{steer} the search --- to
%     decide what to do next, rather than to score what we already did.

%     Concretely, on this URL: flipping \Code{'h'} to \Code{'H'} is valid
%     but covers nothing new, while deleting the \Code{':'} is
%     \emph{invalid} and yet reaches four locations the seed never did
%     --- the error-handling code. Both facts matter: a fuzzer guided by
%     coverage will keep inputs that a fuzzer guided by validity would
%     throw away.
%   \end{solution}

\end{parts}

\end{questions}

\end{document}
